 \documentclass{beamer}
 \usetheme{Warsaw}
 \author{Bert}
 \title{A tale of two primes}
 \begin{document}

 \begin{frame}
 \titlepage
 \end{frame}
 \begin{frame}{Article}
 Some text about the article.
 \end{frame}
 \begin{frame}{Mathematica}
 A helpful tool for mathematicians.
 \end{frame}

 \begin{frame}{Article}
 \begin{block}{Example}
 This is an example of a block.
 \end{block}

\pause

 \begin{block}{Euclid's theorem}
 This is a theorem.
 \end{block}
 \end{frame}

\begin{frame}
\frametitle{Sets}

A \alert{set} is a collection of objects.\uncover<2->{ For example:

\[
Z=\{\text{cow},\text{pig},\text{elephant}\}.
\]
}

\uncover<3->{We call the objects in $Z$ the \alert{elements} of $Z$.}
\uncover<4->{ We write
\[
\text{cow} \in Z
\]
}

\uncover<5->{with ``cow is an element of $Z$''.}
\uncover<6->{ Frequently encountered sets are}

\[
\begin{split}
\uncover<7->{\mathbb{N}} \uncover<8->{= \{1,2,3,\ldots \}}&
\uncover<9->{ \qquad (\text{``natural numbers''})\\}
\uncover<10->{\mathbb{Z}} \uncover<11->{= \{\ldots,-2,-1,0,1,2,\ldots \}}&
\uncover<12->{ \qquad (\text{``integer numbers''})\\}
\uncover<13->{\mathbb{Q}} \uncover<14->{= \{p/q : p,q\in\mathbb{Z} \text{ and } q\neq 0\}}&
\uncover<15->{ \qquad (\text{``rational numbers''})\\}
\uncover<16->{\mathbb{R}} \uncover<17->{= \{\hbox{decimal numbers}\}}&
\uncover<18->{ \qquad (\text{``real numbers''})}
\end{split}
\]

\end{frame}

\begin{frame}
  The derivative of \(f(x) = g(x) \cdot h(x)\), with \(g(x) = x^2\) and \(h(x) = \sin(x)\) equals

  \begin{align*}
    f'(x)
      &\uncover<2->{= g'(x) \cdot h(x)}
       \uncover<3->{+ g(x) \cdot h'(x)} \\
      &\uncover<4->{= 2x \cdot \sin(x)}
       \uncover<5->{+ x^2 \cdot \cos(x).}
  \end{align*}
\end{frame}

\begin{frame}
 \begin{itemize}
 \item<4-6> $[a,b]\uncover<5-6>{=\{x\in\mathbb{R} : a\leq x\leq b\}$,}
 
 \item<6-> $(a,b)\uncover<7->{=\{x\in\mathbb{R} : a< x< b\}$,}

 \item<8-> $(a,\infty)\uncover<9->{=\{x\in\mathbb{R} : x>a\}$.}

 \end{itemize}
\end{frame}

 \end{document}